% OpenStax Calculus Volume 3 -- Section 2.4 Exercises: The Cross Product
% Exercises reproduced from OpenStax, licensed under CC BY 4.0.
% Source: https://openstax.org/books/calculus-volume-3/pages/2-4-the-cross-product
\documentclass{exercises}
\title{OpenStax Calculus Volume 3}
\subtitle{Section 2.4 Exercises: The Cross Product}
\sourceurl{https://openstax.org/books/calculus-volume-3/pages/2-4-the-cross-product}
\author{}
\date{}

\begin{document}
\maketitle


For the following exercises, the vectors $\mathbf{u}$ and $\mathbf{v}$ are given.


(a) Find the cross product $\mathbf{u}\times \mathbf{v}$ of the vectors $\mathbf{u}$ and $\mathbf{v}$. Express the answer in component form.


(b) Sketch the vectors $\mathbf{u},\mathbf{v}$, and $\mathbf{u}\times \mathbf{v}$.

\begin{enumerate}[start=1]
\item $\mathbf{u}=\langle2,0,0\rangle$, $\mathbf{v}=\langle2,2,0\rangle$
\item $\mathbf{u}=\langle3,2,-1\rangle$, $\mathbf{v}=\langle1,1,0\rangle$
\item $\mathbf{u}=2\mathbf{i}+3\mathbf{j}$, $\mathbf{v}=\mathbf{j}+2\mathbf{k}$
\item $\mathbf{u}=2\mathbf{j}+3\mathbf{k}$, $\mathbf{v}=3\mathbf{i}+\mathbf{k}$
\item Simplify $(\mathbf{i}\times \mathbf{i}-2\mathbf{i}\times \mathbf{j}-4\mathbf{i}\times \mathbf{k}+3\mathbf{j}\times \mathbf{k})\times \mathbf{i}$.
\item Simplify $\mathbf{j}\times (\mathbf{k}\times \mathbf{j}+2\mathbf{j}\times \mathbf{i}-3\mathbf{j}\times \mathbf{j}+5\mathbf{i}\times \mathbf{k})$.
\end{enumerate}

In the following exercises, vectors $\mathbf{u}$ and $\mathbf{v}$ are given. Find unit vector $\mathbf{w}$ in the direction of the cross product vector $\mathbf{u}\times \mathbf{v}$. Express your answer using standard unit vectors.

\begin{enumerate}[resume]
\item $\mathbf{u}=\langle3,-1,2\rangle$, $\mathbf{v}=\langle-2,0,1\rangle$
\item $\mathbf{u}=\langle2,6,1\rangle$, $\mathbf{v}=\langle3,0,1\rangle$
\item $\mathbf{u}=\overset{\to}{AB}$, $\mathbf{v}=\overset{\to}{AC}$, where $A(1,0,1)$, $B(1,-1,3)$, and $C(0,0,5)$
\item $\mathbf{u}=\overset{\to}{OP}$, $\mathbf{v}=\overset{\to}{PQ}$, where $P(-1,1,0)$ and $Q(0,2,1)$
\item Determine the real number $\alpha$ such that $\mathbf{u}\times \mathbf{v}$ and $\mathbf{i}$ are orthogonal, where $\mathbf{u}=3\mathbf{i}+\mathbf{j}-5\mathbf{k}$ and $\mathbf{v}=4\mathbf{i}-2\mathbf{j}+\alpha \mathbf{k}$.
\item Show that $\mathbf{u}\times \mathbf{v}$ and $2\mathbf{i}-14\mathbf{j}+2\mathbf{k}$ cannot be orthogonal for any $\alpha$ real number, where $\mathbf{u}=\mathbf{i}+7\mathbf{j}-\mathbf{k}$ and $\mathbf{v}=\alpha \mathbf{i}+5\mathbf{j}+\mathbf{k}$.
\item Show that $\mathbf{u}\times \mathbf{v}$ is orthogonal to $\mathbf{u}+\mathbf{v}$ and $\mathbf{u}-\mathbf{v}$, where $\mathbf{u}$ and $\mathbf{v}$ are nonzero vectors.
\item Show that $\mathbf{v}\times \mathbf{u}$ is orthogonal to $(\mathbf{u}\cdot \mathbf{v})(\mathbf{u}+\mathbf{v})+\mathbf{u}$, where $\mathbf{u}$ and $\mathbf{v}$ are nonzero vectors.
\item Calculate the determinant $\begin{vmatrix}\mathbf{i} & \mathbf{j} & \mathbf{k} \\ 1 & -1 & 7 \\ 2 & 0 & 3\end{vmatrix}$.
\item Calculate the determinant $\begin{vmatrix}\mathbf{i} & \mathbf{j} & \mathbf{k} \\ 0 & 3 & -4 \\ 1 & 6 & -1\end{vmatrix}$.
\end{enumerate}

For the following exercises, the vectors $\mathbf{u}$ and $\mathbf{v}$ are given. Use determinant notation to find vector $\mathbf{w}$ orthogonal to vectors $\mathbf{u}$ and $\mathbf{v}$.

\begin{enumerate}[resume]
\item $\mathbf{u}=\langle-1,0,e^{t}\rangle$, $\mathbf{v}=\langle1,e^{-t},0\rangle$, where $t$ is a real number
\item $\mathbf{u}=\langle1,0,x\rangle$, $\mathbf{v}=\langle \frac{2}{x},1,0\rangle$, where $x$ is a nonzero real number
\item Find vector $(\mathbf{a}-2\mathbf{b})\times \mathbf{c}$, where $\mathbf{a}=\begin{vmatrix}\mathbf{i} & \mathbf{j} & \mathbf{k} \\ 2 & -1 & 5 \\ 0 & 1 & 8\end{vmatrix}$, $\mathbf{b}=\begin{vmatrix}\mathbf{i} & \mathbf{j} & \mathbf{k} \\ 0 & 1 & 1 \\ 2 & -1 & -2\end{vmatrix}$, and $\mathbf{c}=\mathbf{i}+\mathbf{j}+\mathbf{k}$.
\item Find vector $\mathbf{c}\times (\mathbf{a}+3\mathbf{b})$, where $\mathbf{a}=\begin{vmatrix}\mathbf{i} & \mathbf{j} & \mathbf{k} \\ 5 & 0 & 9 \\ 0 & 1 & 0\end{vmatrix}$, $\mathbf{b}=\begin{vmatrix}\mathbf{i} & \mathbf{j} & \mathbf{k} \\ 0 & -1 & 1 \\ 7 & 1 & -1\end{vmatrix}$, and $\mathbf{c}=\mathbf{i}-\mathbf{k}$.
\item \techrequired Use the cross product $\mathbf{u}\times \mathbf{v}$ to find the acute angle between vectors $\mathbf{u}$ and $\mathbf{v}$, where $\mathbf{u}=\mathbf{i}+2\mathbf{j}$ and $\mathbf{v}=\mathbf{i}+\mathbf{k}$. Express the answer in degrees rounded to the nearest integer.
\item \techrequired Use the cross product $\mathbf{u}\times \mathbf{v}$ to find the obtuse angle between vectors $\mathbf{u}$ and $\mathbf{v}$, where $\mathbf{u}=-\mathbf{i}+3\mathbf{j}+\mathbf{k}$ and $\mathbf{v}=\mathbf{i}-2\mathbf{j}$. Express the answer in degrees rounded to the nearest integer.
\item Use the sine and cosine of the angle between two nonzero vectors $\mathbf{u}$ and $\mathbf{v}$ to prove Lagrange's identity: ${\|\mathbf{u}\times \mathbf{v}\|}^{2}={\|\mathbf{u}\|}^{2}{\|\mathbf{v}\|}^{2}-{(\mathbf{u}\cdot \mathbf{v})}^{2}$.
\item Verify Lagrange's identity ${\|\mathbf{u}\times \mathbf{v}\|}^{2}={\|\mathbf{u}\|}^{2}{\|\mathbf{v}\|}^{2}-{(\mathbf{u}\cdot \mathbf{v})}^{2}$ for vectors $\mathbf{u}=-\mathbf{i}+\mathbf{j}-2\mathbf{k}$ and $\mathbf{v}=2\mathbf{i}-\mathbf{j}$.
\item Nonzero vectors $\mathbf{u}$ and $\mathbf{v}$ are called collinear if there exists a nonzero scalar $\alpha$ such that $\mathbf{v}=\alpha \mathbf{u}$. Show that $\mathbf{u}$ and $\mathbf{v}$ are collinear if and only if $\mathbf{u}\times \mathbf{v}=0$.
\item Nonzero vectors $\mathbf{u}$ and $\mathbf{v}$ are called collinear if there exists a nonzero scalar $\alpha$ such that $\mathbf{v}=\alpha \mathbf{u}$. Show that vectors $\overset{\to}{AB}$ and $\overset{\to}{AC}$ are collinear, where $A(4,1,0)$, $B(6,5,-2)$, and $C(5,3,-1)$.
\item Find the area of the parallelogram with adjacent sides $\mathbf{u}=\langle3,2,0\rangle$ and $\mathbf{v}=\langle0,2,1\rangle$.
\item Find the area of the parallelogram with adjacent sides $\mathbf{u}=\mathbf{i}+\mathbf{j}$ and $\mathbf{v}=\mathbf{i}+\mathbf{k}$.
\item Consider points $A(3,-1,2),B(2,1,5)$, and $C(1,-2,-2)$.
\begin{enumerate}[label=(\alph*)]
  \item Find the area of parallelogram $ABCD$ with adjacent sides $\overset{\to}{AB}$ and $\overset{\to}{AC}$.
  \item Find the area of triangle $ABC$.
  \item Find the distance from point $A$ to line $BC$.
\end{enumerate}
\item Consider points $A(2,-3,4),B(0,1,2)$, and $C(-1,2,0)$.
\begin{enumerate}[label=(\alph*)]
  \item Find the area of parallelogram $ABCD$ with adjacent sides $\overset{\to}{AB}$ and $\overset{\to}{AC}$.
  \item Find the area of triangle $ABC$.
  \item Find the distance from point $B$ to line $AC$.
\end{enumerate}
\end{enumerate}

In the following exercises, vectors $\mathbf{u}$, $\mathbf{v}$, and $\mathbf{w}$ are given.


(a) Find the triple scalar product $\mathbf{u}\cdot(\mathbf{v}\times \mathbf{w})$.


(b) Find the volume of the parallelepiped with the adjacent edges $\mathbf{u}$, $\mathbf{v}$, and $\mathbf{w}$.

\begin{enumerate}[resume]
\item $\mathbf{u}=\mathbf{i}+\mathbf{j}$, $\mathbf{v}=\mathbf{j}+\mathbf{k}$, and $\mathbf{w}=\mathbf{i}+\mathbf{k}$
\item $\mathbf{u}=\langle-3,5,-1\rangle$, $\mathbf{v}=\langle0,2,-2\rangle$, and $\mathbf{w}=\langle3,1,1\rangle$
\item Calculate the triple scalar products $\mathbf{v}\cdot(\mathbf{u}\times \mathbf{w})$ and $\mathbf{w}\cdot(\mathbf{u}\times \mathbf{v})$, where $\mathbf{u}=\langle1,1,1\rangle$, $\mathbf{v}=\langle7,6,9\rangle$, and $\mathbf{w}=\langle4,2,7\rangle$.
\item Calculate the triple scalar products $\mathbf{w}\cdot(\mathbf{v}\times \mathbf{u})$ and $\mathbf{u}\cdot(\mathbf{w}\times \mathbf{v})$, where $\mathbf{u}=\langle4,2,-1\rangle$, $\mathbf{v}=\langle2,5,-3\rangle$, and $\mathbf{w}=\langle9,5,-10\rangle$.
\item Find vectors $\mathbf{a}$, $\mathbf{b}$, and $\mathbf{c}$ with a triple scalar product given by the determinant $\begin{vmatrix}1 & 2 & 3 \\ 0 & 2 & 5 \\ 8 & 9 & 2\end{vmatrix}$. Determine their triple scalar product.
\item The triple scalar product of vectors $\mathbf{a}$, $\mathbf{b}$, and $\mathbf{c}$ is given by the determinant $\begin{vmatrix}0 & -2 & 1 \\ 0 & 1 & 4 \\ 1 & -3 & 7\end{vmatrix}$. Find vector $\mathbf{a}-\mathbf{b}+\mathbf{c}$.
\item Consider the parallelepiped with edges $OA,OB$, and $OC$, where $A(2,1,0),B(1,2,0)$, and $C(0,1,\alpha)$.
\begin{enumerate}[label=(\alph*)]
  \item Find the real number $\alpha>0$ such that the volume of the parallelepiped is $3$ units$^{3}$.
  \item For $\alpha=1$, find the height $h$ from vertex $C$ of the parallelepiped to the plane formed by the edges $OA$ and $OB$.
\end{enumerate}
\item Consider points $A(\alpha,0,0),B(0,\beta,0)$, and $C(0,0,\gamma)$, with $\alpha$, $\beta$, and $\gamma$ positive real numbers.
\begin{enumerate}[label=(\alph*)]
  \item Determine the volume of the parallelepiped with adjacent sides $\overset{\to}{OA}$, $\overset{\to}{OB}$, and $\overset{\to}{OC}$.
  \item Find the volume of the tetrahedron with vertices $O$, $A$, $B$, and $C$. (Hint: The volume of the tetrahedron is $1/6$ of the volume of the parallelepiped.)
  \item Find the distance from the origin to the plane determined by $A$, $B$, and $C$. Sketch the parallelepiped and tetrahedron.
\end{enumerate}
\item Let $\mathbf{u}$, $\mathbf{v}$, and $\mathbf{w}$ be three-dimensional vectors and $c$ be a real number. Prove the following properties of the cross product.
\begin{enumerate}[label=(\alph*)]
  \item $\mathbf{u}\times \mathbf{u}=0$
  \item $\mathbf{u}\times (\mathbf{v}+\mathbf{w})=(\mathbf{u}\times \mathbf{v})+(\mathbf{u}\times \mathbf{w})$
  \item $c(\mathbf{u}\times \mathbf{v})=(c\mathbf{u})\times \mathbf{v}=\mathbf{u}\times (c\mathbf{v})$
  \item $\mathbf{u}\cdot(\mathbf{u}\times \mathbf{v})=0$
\end{enumerate}
\item Show that vectors $\mathbf{u}=\langle1,0,-8\rangle$, $\mathbf{v}=\langle0,1,6\rangle$, and $\mathbf{w}=\langle-1,9,3\rangle$ satisfy the following properties of the cross product.
\begin{enumerate}[label=(\alph*)]
  \item $\mathbf{u}\times \mathbf{u}=0$
  \item $\mathbf{u}\times (\mathbf{v}+\mathbf{w})=(\mathbf{u}\times \mathbf{v})+(\mathbf{u}\times \mathbf{w})$
  \item $c(\mathbf{u}\times \mathbf{v})=(c\mathbf{u})\times \mathbf{v}=\mathbf{u}\times (c\mathbf{v})$
  \item $\mathbf{u}\cdot(\mathbf{u}\times \mathbf{v})=0$
\end{enumerate}
\item Nonzero vectors $\mathbf{u}$, $\mathbf{v}$, and $\mathbf{w}$ are said to be linearly dependent if one of the vectors is a linear combination of the other two. For instance, there exist two nonzero real numbers $\alpha$ and $\beta$ such that $\mathbf{w}=\alpha \mathbf{u}+\beta \mathbf{v}$. Otherwise, the vectors are called linearly independent. Show that $\mathbf{u}$, $\mathbf{v}$, and $\mathbf{w}$ are coplanar if and only if they are linear dependent.
\item Consider vectors $\mathbf{u}=\langle1,4,-7\rangle$, $\mathbf{v}=\langle2,-1,4\rangle$, $\mathbf{w}=\langle0,-9,18\rangle$, and $\mathbf{p}=\langle0,-9,17\rangle$.
\begin{enumerate}[label=(\alph*)]
  \item Show that $\mathbf{u}$, $\mathbf{v}$, and $\mathbf{w}$ are coplanar by using their triple scalar product
  \item Show that $\mathbf{u}$, $\mathbf{v}$, and $\mathbf{w}$ are coplanar, using the definition that there exist two nonzero real numbers $\alpha$ and $\beta$ such that $\mathbf{w}=\alpha \mathbf{u}+\beta \mathbf{v}$.
  \item Show that $\mathbf{u}$, $\mathbf{v}$, and $\mathbf{p}$ are linearly independent---that is, none of the vectors is a linear combination of the other two.
\end{enumerate}
\item Consider points $A(0,0,2)$, $B(1,0,2)$, $C(1,1,2)$, and $D(0,1,2)$. Are vectors $\overset{\to}{AB}$, $\overset{\to}{AC}$, and $\overset{\to}{AD}$ linearly dependent (that is, one of the vectors is a linear combination of the other two)?
\item Show that vectors $\mathbf{i}+\mathbf{j}$, $\mathbf{i}-\mathbf{j}$, and $\mathbf{i}+\mathbf{j}+\mathbf{k}$ are linearly independent---that is, there do not exist two nonzero real numbers $\alpha$ and $\beta$ such that $\mathbf{i}+\mathbf{j}+\mathbf{k}=\alpha(\mathbf{i}+\mathbf{j})+\beta(\mathbf{i}-\mathbf{j})$.
\item Let $\mathbf{u}=\langle u_{1},u_{2}\rangle$ and $\mathbf{v}=\langle v_{1},v_{2}\rangle$ be two-dimensional vectors. The cross product of vectors $\mathbf{u}$ and $\mathbf{v}$ is not defined. However, if the vectors are regarded as the three-dimensional vectors $\tilde{\mathbf{u}}=\langle u_{1},u_{2},0\rangle$ and $\tilde{\mathbf{v}}=\langle v_{1},v_{2},0\rangle$, respectively, then, in this case, we can define the cross product of $\tilde{\mathbf{u}}$ and $\tilde{\mathbf{v}}$. In particular, in determinant notation, the cross product of $\tilde{\mathbf{u}}$ and $\tilde{\mathbf{v}}$ is given by \[\tilde{\mathbf{u}}\times \tilde{\mathbf{v}}=\begin{vmatrix}\mathbf{i} & \mathbf{j} & \mathbf{k} \\ u_{1} & u_{2} & 0 \\ v_{1} & v_{2} & 0\end{vmatrix}.\] Use this result to compute $(\mathbf{i}\cos \theta+\mathbf{j}\sin \theta)\times (\mathbf{i}\sin \theta-\mathbf{j}\cos \theta)$, where $\theta$ is a real number.
\item Consider points $P(2,1)$, $Q(4,2)$, and $R(1,2)$.
\begin{enumerate}[label=(\alph*)]
  \item Find the area of triangle $P$, $Q$, and $R$.
  \item Determine the distance from point $R$ to the line passing through $P$ and $Q$.
\end{enumerate}
\item Determine a vector of magnitude $10$ perpendicular to the plane passing through the x-axis and point $P(1,2,4)$.
\item Determine a unit vector perpendicular to the plane passing through the z-axis and point $A(3,1,-2)$.
\item Consider $\mathbf{u}$ and $\mathbf{v}$ two three-dimensional vectors. If the magnitude of the cross product vector $\mathbf{u}\times \mathbf{v}$ is $k$ times larger than the magnitude of vector $\mathbf{u}$, show that the magnitude of $\mathbf{v}$ is greater than or equal to $k$, where $k$ is a natural number.
\item \techrequired Assume that the magnitudes of two nonzero vectors $\mathbf{u}$ and $\mathbf{v}$ are known. The function $f(\theta)=\|\mathbf{u}\|\|\mathbf{v}\|\sin \theta$ defines the magnitude of the cross product vector $\mathbf{u}\times \mathbf{v}$, where $\theta \in[0,\pi]$ is the angle between $\mathbf{u}$ and $\mathbf{v}$.
\begin{enumerate}[label=(\alph*)]
  \item Graph the function $f$.
  \item Find the absolute minimum and maximum of function $f$. Interpret the results.
  \item If $\|\mathbf{u}\|=5$ and $\|\mathbf{v}\|=2$, find the angle between $\mathbf{u}$ and $\mathbf{v}$ if the magnitude of their cross product vector is equal to $9$.
\end{enumerate}
\item Find all vectors $\mathbf{w}=\langle w_{1},w_{2},w_{3}\rangle$ that satisfy the equation $\langle1,1,1\rangle \times \mathbf{w}=\langle-1,-1,2\rangle$.
\item Solve the equation $\mathbf{w}\times \langle1,0,-1\rangle=\langle3,0,3\rangle$, where $\mathbf{w}=\langle w_{1},w_{2},w_{3}\rangle$ is a nonzero vector with a magnitude of $3$.
\item \techrequired A mechanic uses a 12-in. wrench to turn a bolt. The wrench makes a $30^\circ$ angle with the horizontal. If the mechanic applies a vertical force of $10$ lb on the wrench handle, what is the magnitude of the torque at point $P$ (see the following figure)? Express the answer in foot-pounds rounded to two decimal places. \\ \textit{[Figure: This figure is the image of an open-end wrench. The lower portion of the wrench is at point P. The wrench has a length of ``12 in.'' The angle the wrench makes with a horizontal line from P is 30 degrees. At the top of the wrench is a downward vertical vector labeled ``10 lb.'']}
\item \techrequired A boy applies the brakes on a bicycle by applying a downward force of $20$ lb on the pedal when the 6-in. crank makes a $40^\circ$ angle with the horizontal (see the following figure). Find the torque at point $P$. Express your answer in foot-pounds rounded to two decimal places. \\ \textit{[Figure: This figure shows the pedals, cranks, and chain of a bicycle. The distance along the crank to the top pedal is 6 in. The angle of the crank is 40 degrees with the horizontal, measured toward the rear. The top pedal has a downward vector labeled ``20 lb''.]}
\item \techrequired Find the magnitude of the force that needs to be applied to the end of a 20-cm wrench located on the positive direction of the y-axis if the force is applied in the direction $\langle0,1,-2\rangle$ and it produces a $100$ N·m torque to the bolt located at the origin.
\item \techrequired What is the magnitude of the force required to be applied to the end of a 1-ft wrench at an angle of $35^\circ$ to produce a torque of $20$ ft-lbs?
\item \techrequired The force vector $\mathbf{F}$ acting on a proton with an electric charge of $1.6\times 10^{-19}\text{C}$ (in coulombs) moving in a magnetic field $\mathbf{B}$ where the velocity vector $\mathbf{v}$ is given by $\mathbf{F}=1.6\times 10^{-19}(\mathbf{v}\times \mathbf{B})$ (here, $\mathbf{v}$ is expressed in meters per second, $\mathbf{B}$ is in tesla [T], and $\mathbf{F}$ is in newtons [N]). Find the force that acts on a proton that moves in the xy-plane at velocity $\mathbf{v}=10^{5}\mathbf{i}+10^{5}\mathbf{j}$ (in meters per second) in a magnetic field given by $\mathbf{B}=0.3\mathbf{j}$.
\item \techrequired The force vector $\mathbf{F}$ acting on a proton with an electric charge of $1.6\times 10^{-19}\text{C}$ moving in a magnetic field $\mathbf{B}$ where the velocity vector $\mathbf{v}$ is given by $\mathbf{F}=1.6\times 10^{-19}(\mathbf{v}\times \mathbf{B})$ (here, $\mathbf{v}$ is expressed in meters per second, $\mathbf{B}$ in $\text{T}$, and $\mathbf{F}$ in $\text{N}$). If the magnitude of force $\mathbf{F}$ acting on a proton is $5.9\times 10^{-17}$ N and the proton is moving at the speed of 300 m/sec in magnetic field $\mathbf{B}$ of magnitude 2.4 T, find the angle between velocity vector $\mathbf{v}$ of the proton and magnetic field $\mathbf{B}$. Express the answer in degrees rounded to the nearest integer.
\item \techrequired Consider $\mathbf{r}(t)=\langle \cos t,\sin t,2t\rangle$ the position vector of a particle at time $t\in[0,30]$, where the components of $\mathbf{r}$ are expressed in centimeters and time in seconds. Let $\overset{\to}{OP}$ be the position vector of the particle after $1$ sec.
\begin{enumerate}[label=(\alph*)]
  \item Determine unit vector $\mathbf{B}(t)$ (called the binormal unit vector) that has the direction of cross product vector $\mathbf{v}(t)\times \mathbf{a}(t)$, where $\mathbf{v}(t)$ and $\mathbf{a}(t)$ are the instantaneous velocity vector and, respectively, the acceleration vector of the particle after $t$ seconds.
  \item Use a CAS to visualize vectors $\mathbf{v}(1)$, $\mathbf{a}(1)$, and $\mathbf{B}(1)$ as vectors starting at point $P$ along with the path of the particle.
\end{enumerate}
\item A solar panel is mounted on the roof of a house. The panel may be regarded as positioned at the points of coordinates (in meters) $A(8,0,0)$, $B(8,18,0)$, $C(0,18,8)$, and $D(0,0,8)$ (see the following figure). \\ \textit{[Figure: This figure shows a rectangular set of solar panels on a roof. The corners are labeled ``A, B, C, D.'' Also there is a vector drawn from A to D. There is another vector along the bottom of the rectangle from A to B.]}
\begin{enumerate}[label=(\alph*)]
  \item Find vector $\mathbf{n}=\overset{\to}{AB}\times \overset{\to}{AD}$ perpendicular to the surface of the solar panels. Express the answer using standard unit vectors.
  \item Assume unit vector $\mathbf{s}=\frac{1}{\sqrt{3}}\mathbf{i}+\frac{1}{\sqrt{3}}\mathbf{j}+\frac{1}{\sqrt{3}}\mathbf{k}$ points toward the Sun at a particular time of the day and the flow of solar energy is $\mathbf{F}=900\mathbf{s}$ (in watts per square meter [$\text{W/m}^{2}$]). Find the predicted amount of electrical power the panel can produce, which is given by the dot product of vectors $\mathbf{F}$ and $\mathbf{n}$ (expressed in watts).
  \item Determine the angle of elevation of the Sun above the solar panel. Express the answer in degrees rounded to the nearest whole number. (Hint: The angle between vectors $\mathbf{n}$ and $\mathbf{s}$ and the angle of elevation are complementary.)
\end{enumerate}
\end{enumerate}

\end{document}
